% ---- pg-379 （书页 367） ----

\noindent 8.11\quad 考虑\hyperref[tab:8.5]{表 8.5} 的人口普查数据. 对 $X_5=$ 家庭收入中位数的观察值, 用千美元计数而不用万美元. 也就是说, 将表中第 6 列的所有数字都乘以 10.

\begin{mvtable}
\phantomsection\label{tab:8.5}\mvtabcap{表 8.5\quad 人口调查数据}
\mvfit{%
\begin{tabular}{lccccc}
\toprule
区 & 总人口(千人) & 职业度(\%) & 16 岁以上就业率(\%) & 政府雇用率(\%) & 家庭收入中位数(10 万美元)\\
\midrule
1 & 2.67 & 5.71 & 69.02 & 30.3 & 1.48\\
2 & 2.25 & 4.37 & 72.98 & 43.3 & 1.44\\
3 & 3.12 & 10.27 & 64.94 & 32.0 & 2.11\\
4 & 5.14 & 7.44 & 71.29 & 24.5 & 1.85\\
5 & 5.54 & 9.25 & 74.94 & 31.0 & 2.23\\
6 & 5.04 & 4.84 & 53.61 & 48.2 & 1.60\\
7 & 3.14 & 4.82 & 67.00 & 37.6 & 1.52\\
8 & 2.43 & 2.40 & 67.20 & 36.8 & 1.40\\
9 & 5.38 & 4.30 & 83.03 & 19.7 & 2.07\\
10 & 7.34 & 2.73 & 72.60 & 24.5 & 1.42\\
$\vdots$ & $\vdots$ & $\vdots$ & $\vdots$ & $\vdots$ & $\vdots$\\
52 & 7.25 & 1.16 & 78.52 & 23.6 & 1.50\\
53 & 5.44 & 2.93 & 73.59 & 22.3 & 1.65\\
54 & 5.83 & 4.47 & 77.33 & 26.2 & 2.16\\
55 & 3.74 & 2.26 & 79.70 & 20.2 & 1.58\\
56 & 9.21 & 2.36 & 74.58 & 21.8 & 1.72\\
57 & 2.14 & 6.30 & 86.54 & 17.4 & 2.80\\
58 & 6.62 & 4.79 & 78.84 & 20.0 & 2.33\\
59 & 4.24 & 5.82 & 71.39 & 27.1 & 1.69\\
60 & 4.72 & 4.71 & 78.01 & 20.6 & 1.55\\
61 & 6.48 & 4.93 & 74.23 & 20.9 & 1.98\\
\bottomrule
\end{tabular}
}

{\small 注: 观测数据似乎有相关性, 即 61 个观测可能不构成随机样本. 全部数据可从 www.prenhall.com/statistics 得到.}
\end{mvtable}

(a) 构造人口普查数据的样本协方差矩阵 $\mathbf{S}$, 此时 $X_5=$ 家庭收入中位数用千美元为单位计数. (注意, 这个协方差矩阵可由例 8.3 给出的协方差矩阵得到, 将其第 5 行和第 5 列的非对角元乘以 10 并将对角元 $s_{55}$ 乘以 100 即可, 为什么呢?)

(b) 求特征值-特征向量对及(a)中协方差矩阵的前两个样本主成分.

(c) 计算总方差中用(b)求得的前两个主成分解释的比例, 求相关系数 $r_{\hat{y}_i,x_k}$. 如果可能的话, 解释这些成分, 并将所得结果与例 8.3 的结果进行比较. 计数尺度上的变化对主成分有什么影响?

\noindent 8.12\quad 考虑列于\hyperref[tab:1.5]{表 1.5} 的空气污染数据, 如果可能请用小于 $p=7$ 维来综合这些数据. 用协方差矩阵 $\mathbf{S}$ 和相关系数矩阵 $\mathbf{R}$ 对这些数据做主成分分析, 你能从中学到什么? 选择哪个矩阵作分析有不同吗? 这些数据可由三维甚至更少的维数来综合吗? 你能解释这些主成分吗?

\noindent 8.13\quad 放射疗法数据在\hyperref[tab:1.7]{表 1.7} 中, 对 $p=6$ 个代表病人对放射疗法的反应的变量有 $n=98$ 次观察值.

(a) 求这些数据的协方差矩阵 $\mathbf{S}$ 及相关系数矩阵 $\mathbf{R}$.

(b) 挑选一个矩阵 $\mathbf{S}$ 或 $\mathbf{R}$(证明你的选择合理)然后求特征值和特征向量. 按从大到小的
